\documentclass[10pt,a4paper]{amsart}
\usepackage{amsmath,amssymb,amsthm,graphicx}
% Retained minimal preamble for the Stacks Project chapter sites-cohomology.tex (Cohomology on Sites).
% Verbatim excerpt of the chapter preamble: Xy-pic setup (native preamble lines 14-18) and the native
% theorem declarations with the mathematical macros (native preamble lines 154-217). Copied mechanically;
% nothing added, removed or redefined. The wrapper class and the amsmath/amssymb/amsthm/graphicx base
% packages are supplied by the assembler, so the class/externaldocument part of the native preamble is omitted.
% For commutative diagrams we use Xy-pic
\usepackage[all]{xy}

% We use 2cell for 2-commutative diagrams.
\xyoption{2cell}
% Theorem environments.
%
\theoremstyle{plain}
\newtheorem{theorem}[subsection]{Theorem}
\newtheorem{proposition}[subsection]{Proposition}
\newtheorem{lemma}[subsection]{Lemma}

\theoremstyle{definition}
\newtheorem{definition}[subsection]{Definition}
\newtheorem{example}[subsection]{Example}
\newtheorem{exercise}[subsection]{Exercise}
\newtheorem{situation}[subsection]{Situation}

\theoremstyle{remark}
\newtheorem{remark}[subsection]{Remark}
\newtheorem{remarks}[subsection]{Remarks}

\numberwithin{equation}{subsection}

% Macros
%
\def\lim{\mathop{\mathrm{lim}}\nolimits}
\def\colim{\mathop{\mathrm{colim}}\nolimits}
\def\Spec{\mathop{\mathrm{Spec}}}
\def\Hom{\mathop{\mathrm{Hom}}\nolimits}
\def\Ext{\mathop{\mathrm{Ext}}\nolimits}
\def\SheafHom{\mathop{\mathcal{H}\!\mathit{om}}\nolimits}
\def\SheafExt{\mathop{\mathcal{E}\!\mathit{xt}}\nolimits}
\def\Sch{\mathit{Sch}}
\def\Mor{\mathop{\mathrm{Mor}}\nolimits}
\def\Ob{\mathop{\mathrm{Ob}}\nolimits}
\def\Sh{\mathop{\mathit{Sh}}\nolimits}
\def\NL{\mathop{N\!L}\nolimits}
\def\CH{\mathop{\mathrm{CH}}\nolimits}
\def\proetale{{pro\text{-}\acute{e}tale}}
\def\etale{{\acute{e}tale}}
\def\QCoh{\mathit{QCoh}}
\def\Ker{\mathop{\mathrm{Ker}}}
\def\Im{\mathop{\mathrm{Im}}}
\def\Coker{\mathop{\mathrm{Coker}}}
\def\Coim{\mathop{\mathrm{Coim}}}

% Boxtimes
%
\DeclareMathSymbol{\boxtimes}{\mathbin}{AMSa}{"02}

%
% Macros for moduli stacks/spaces
%
\def\QCohstack{\mathcal{QC}\!\mathit{oh}}
\def\Cohstack{\mathcal{C}\!\mathit{oh}}
\def\Spacesstack{\mathcal{S}\!\mathit{paces}}
\def\Quotfunctor{\mathrm{Quot}}
\def\Hilbfunctor{\mathrm{Hilb}}
\def\Curvesstack{\mathcal{C}\!\mathit{urves}}
\def\Polarizedstack{\mathcal{P}\!\mathit{olarized}}
\def\Complexesstack{\mathcal{C}\!\mathit{omplexes}}
% \Pic is the operator that assigns to X its picard group, usage \Pic(X)
% \Picardstack_{X/B} denotes the Picard stack of X over B
% \Picardfunctor_{X/B} denotes the Picard functor of X over B
\def\Pic{\mathop{\mathrm{Pic}}\nolimits}
\def\Picardstack{\mathcal{P}\!\mathit{ic}}
\def\Picardfunctor{\mathrm{Pic}}
\def\Deformationcategory{\mathcal{D}\!\mathit{ef}}

% --- appended 2026-10-12 (Chief): the project own definitions of the reference/slogan/history
% display environments, copied mechanically from the shared Stacks preamble lines 8-12
% (\usepackage{verbatim} plus the three \newenvironment lines, which the chapter preamble keeps
% verbatim from the project preamble). Nothing added, removed or redefined: these environments
% carry no mathematical content and the project itself typesets them as comments. Without these
% three lines a counted statement whose own text uses them fails with Environment slogan undefined.
\usepackage{verbatim}
\newenvironment{reference}{\comment}{\endcomment}
\newenvironment{slogan}{\comment}{\endcomment}
\newenvironment{history}{\comment}{\endcomment}
\usepackage{cleveref}
\crefname{originallemma}{original lemma}{original lemmas}
\Crefname{originallemma}{Original lemma}{Original lemmas}
\makeatletter
\newlabel{lemma-unbounded-cohomology-of-open}{{0D6F}{source}{}{sourcepointer-1}{}}
\newlabel{lemma-unbounded-cohomology-of-open@cref}{{[originallemma][0][]0D6F}{[1][0][]source}{}{}{}}
\makeatother
\crefname{originallemma}{original lemma}{original lemmas}
\Crefname{originallemma}{Original lemma}{Original lemmas}
\makeatletter
\newlabel{lemma-unbounded-describe-higher-direct-images}{{0D6I}{source}{}{sourcepointer-2}{}}
\newlabel{lemma-unbounded-describe-higher-direct-images@cref}{{[originallemma][0][]0D6I}{[1][0][]source}{}{}{}}
\makeatother
\crefname{originallemma}{original lemma}{original lemmas}
\Crefname{originallemma}{Original lemma}{Original lemmas}
\makeatletter
\newlabel{sites-modules-lemma-module-quotient-flat}{{03EW}{source}{}{sourcepointer-3}{}}
\newlabel{sites-modules-lemma-module-quotient-flat@cref}{{[originallemma][0][]03EW}{[1][0][]source}{}{}{}}
\makeatother
\crefname{originallemma}{original lemma}{original lemmas}
\Crefname{originallemma}{Original lemma}{Original lemmas}
\makeatletter
\newlabel{derived-lemma-subcategory-left-resolution}{{05T7}{source}{}{sourcepointer-4}{}}
\newlabel{derived-lemma-subcategory-left-resolution@cref}{{[originallemma][0][]05T7}{[1][0][]source}{}{}{}}
\makeatother
\crefname{originalsection}{original section}{original sections}
\Crefname{originalsection}{Original section}{Original sections}
\makeatletter
\newlabel{section-flat}{{06YL}{source}{}{sourcepointer-5}{}}
\newlabel{section-flat@cref}{{[originalsection][0][]06YL}{[1][0][]source}{}{}{}}
\makeatother
\crefname{originallemma}{original lemma}{original lemmas}
\Crefname{originallemma}{Original lemma}{Original lemmas}
\makeatletter
\newlabel{lemma-derived-tor-quasi-isomorphism-other-side}{{06YT}{source}{}{sourcepointer-6}{}}
\newlabel{lemma-derived-tor-quasi-isomorphism-other-side@cref}{{[originallemma][0][]06YT}{[1][0][]source}{}{}{}}
\makeatother
\crefname{originallemma}{original lemma}{original lemmas}
\Crefname{originallemma}{Original lemma}{Original lemmas}
\makeatletter
\newlabel{lemma-bounded-flat-K-flat}{{06YQ}{source}{}{sourcepointer-7}{}}
\newlabel{lemma-bounded-flat-K-flat@cref}{{[originallemma][0][]06YQ}{[1][0][]source}{}{}{}}
\makeatother
\title{Cohomology of a complex of modules agrees with that of its underlying complex of abelian sheaves, unboundedly}
\author{The Stacks Project Authors}
\date{stacks-project:sites-cohomology.tex:local-snapshot-20261011; GNU FDL 1.2 or later}
\begin{document}
\maketitle
\noindent\textit{Editorial note.} All delivered mathematical text is the author's own, copied from the pinned Stacks Project chapter file sites-cohomology.tex; no statement or proof step was rewritten, reordered or removed.\par\medskip
\noindent\textit{Editorial note.} The proof is the author's complete argument. The phrase 'see discussion in Section ...' inside it is a pointer to a section of the source tree and not an omission; the closing sentence 'We omit the detailed verification that this construction is indeed an inverse.' is a routine tail. Both are kept verbatim inside the counted proof span.\par\medskip
\noindent\textit{Editorial note.} Every labelled reference occurring inside the counted spans is carried in assembly.external\_references as a non-load-bearing navigation pointer (same-chapter ones at their definition site, cross-chapter ones at the invocation inside the counted span), so no load-bearing mathematics is replaced by a stub. No bibliography citation occurs inside the counted spans.\par\medskip
\noindent\textit{Editorial note.} The checked chapter list for same-name collisions is the set of stacks chapters staged in this repository: chow, cohomology, divisors, etale-cohomology, flat, more-algebra, more-morphisms, morphisms, spaces-more-morphisms, varieties together with this card's chapter sites-cohomology.\par\medskip
\noindent\textit{Editorial note.} The preamble is the chapter's own retained minimal preamble (sources/stacks-sites-cohomology/retained-preamble.tex), a verbatim excerpt of the native preamble: the Xy-pic setup and the native theorem declarations with the mathematical macros.\par\medskip
\noindent\textit{Allowed standard prerequisites (editorial).}\begin{itemize}
\item Unbounded derived category of modules over a ringed site and derived sections and pushforward. Same chapter, the functors RGamma and Rf\_* on the unbounded derived category and their computation from K-injective representatives. Supplies the maps whose being isomorphisms is asserted.
\item K-flat resolutions with flat terms and cofinality of such resolutions. Same chapter, the existence of K-flat resolutions with flat terms and the cofinality of the quasi-isomorphisms from bounded above complexes of flat modules, together with the companion isomorphism between the cohomology of a module and of its underlying abelian sheaf at a point. Supplies the construction of the inverse map in the proof.
\end{itemize}\medskip
\ifdefined\hypertarget\hypertarget{sourcepointer-1}{}\fi
\noindent\textit{Source reference.} Original lemma 0D6F: sources/stacks-sites-cohomology/pinned-original-source.tex, source lines 4329--4329. Navigation pointer invoked at proof line 4485; the counted span contains the pointer only and the cited result is not reproduced inside it, so the pointer is registered as a non-load-bearing navigation stub in the established form for Stacks material. No bibliography citation is involved.\par\medskip
\ifdefined\hypertarget\hypertarget{sourcepointer-2}{}\fi
\noindent\textit{Source reference.} Original lemma 0D6I: sources/stacks-sites-cohomology/pinned-original-source.tex, source lines 4434--4434. Navigation pointer invoked at proof line 4486; the counted span contains the pointer only and the cited result is not reproduced inside it, so the pointer is registered as a non-load-bearing navigation stub in the established form for Stacks material. No bibliography citation is involved.\par\medskip
\ifdefined\hypertarget\hypertarget{sourcepointer-3}{}\fi
\noindent\textit{Source reference.} Original lemma 03EW: sources/stacks-sites-cohomology/pinned-original-source.tex, source lines 4518--4518. Navigation pointer invoked at proof line 4518; the counted span contains the pointer only and the cited result is not reproduced inside it, so the pointer is registered as a non-load-bearing navigation stub in the established form for Stacks material. No bibliography citation is involved.\par\medskip
\ifdefined\hypertarget\hypertarget{sourcepointer-4}{}\fi
\noindent\textit{Source reference.} Original lemma 05T7: sources/stacks-sites-cohomology/pinned-original-source.tex, source lines 4519--4519. Navigation pointer invoked at proof line 4519; the counted span contains the pointer only and the cited result is not reproduced inside it, so the pointer is registered as a non-load-bearing navigation stub in the established form for Stacks material. No bibliography citation is involved.\par\medskip
\ifdefined\hypertarget\hypertarget{sourcepointer-5}{}\fi
\noindent\textit{Source reference.} Original section 06YL: sources/stacks-sites-cohomology/pinned-original-source.tex, source lines 2961--2961. Navigation pointer invoked at proof line 4520; the counted span contains the pointer only and the cited result is not reproduced inside it, so the pointer is registered as a non-load-bearing navigation stub in the established form for Stacks material. No bibliography citation is involved.\par\medskip
\ifdefined\hypertarget\hypertarget{sourcepointer-6}{}\fi
\noindent\textit{Source reference.} Original lemma 06YT: sources/stacks-sites-cohomology/pinned-original-source.tex, source lines 3233--3233. Navigation pointer invoked at proof line 4540; the counted span contains the pointer only and the cited result is not reproduced inside it, so the pointer is registered as a non-load-bearing navigation stub in the established form for Stacks material. No bibliography citation is involved.\par\medskip
\ifdefined\hypertarget\hypertarget{sourcepointer-7}{}\fi
\noindent\textit{Source reference.} Original lemma 06YQ: sources/stacks-sites-cohomology/pinned-original-source.tex, source lines 3120--3120. Navigation pointer invoked at proof line 4541; the counted span contains the pointer only and the cited result is not reproduced inside it, so the pointer is registered as a non-load-bearing navigation stub in the established form for Stacks material. No bibliography citation is involved.\par\medskip
\begin{lemma}
\label{lemma-modules-abelian-unbounded}
Let $(\mathcal{C}, \mathcal{O}_\mathcal{C})$ be a ringed site.
Let $K$ be an object of $D(\mathcal{O}_\mathcal{C})$
and denote $K_{ab}$ its image in $D(\underline{\mathbf{Z}}_\mathcal{C})$.
\begin{enumerate}
\item There is a canonical map
$R\Gamma(\mathcal{C}, K) \to R\Gamma(\mathcal{C}, K_{ab})$
which is an isomorphism in $D(\textit{Ab})$.
\item For any $U \in \mathcal{C}$ there is a canonical map
$R\Gamma(U, K) \to R\Gamma(U, K_{ab})$
which is an isomorphism in $D(\textit{Ab})$.
\item Let $f : (\mathcal{C}, \mathcal{O}_\mathcal{C}) \to
(\mathcal{D}, \mathcal{O}_\mathcal{D})$ be a morphism of ringed sites.
There is a canonical map $Rf_*K \to Rf_*(K_{ab})$ which
is an isomorphism in $D(\underline{\mathbf{Z}}_\mathcal{D})$.
\end{enumerate}
\end{lemma}

\begin{proof}
The map is constructed as follows. Choose a K-injective complex
$\mathcal{I}^\bullet$ representing $K$. Choose a quasi-isomorpism
$\mathcal{I}^\bullet \to \mathcal{J}^\bullet$ where $\mathcal{J}^\bullet$
is a K-injective complex of abelian groups. Then the map in
(1) is given by
$\Gamma(\mathcal{C}, \mathcal{I}^\bullet) \to
\Gamma(\mathcal{C}, \mathcal{J}^\bullet)$
(2) is given by
$\Gamma(U, \mathcal{I}^\bullet) \to \Gamma(U, \mathcal{J}^\bullet)$
and the map in (3) is given by
$f_*\mathcal{I}^\bullet \to f_*\mathcal{J}^\bullet$.
To show that these maps are isomorphisms, it suffices to prove
they induce isomorphisms on cohomology groups and cohomology sheaves.
By Lemmas \ref{lemma-unbounded-cohomology-of-open} and
\ref{lemma-unbounded-describe-higher-direct-images}
it suffices to show that the map
$$
H^0(\mathcal{C}, K) \longrightarrow H^0(\mathcal{C}, K_{ab})
$$
is an isomorphism. Observe that
$$
H^0(\mathcal{C}, K) =
\Hom_{D(\mathcal{O}_\mathcal{C})}(\mathcal{O}_\mathcal{C}, K)
$$
and similarly for the other group. Choose any complex $\mathcal{K}^\bullet$
of $\mathcal{O}_\mathcal{C}$-modules representing $K$. By construction of the
derived category as a localization we have
$$
\Hom_{D(\mathcal{O}_\mathcal{C})}(\mathcal{O}_\mathcal{C}, K) =
\colim_{s : \mathcal{F}^\bullet \to \mathcal{O}_\mathcal{C}}
\Hom_{K(\mathcal{O}_\mathcal{C})}(\mathcal{F}^\bullet, \mathcal{K}^\bullet)
$$
where the colimit is over quasi-isomorphisms $s$ of complexes of
$\mathcal{O}_\mathcal{C}$-modules. Similarly, we have
$$
\Hom_{D(\underline{\mathbf{Z}}_\mathcal{C})}
(\underline{\mathbf{Z}}_\mathcal{C}, K) =
\colim_{s : \mathcal{G}^\bullet \to \underline{\mathbf{Z}}_\mathcal{C}}
\Hom_{K(\underline{\mathbf{Z}}_\mathcal{C})}
(\mathcal{G}^\bullet, \mathcal{K}^\bullet)
$$
Next, we observe that the quasi-isomorphisms
$s : \mathcal{G}^\bullet \to \underline{\mathbf{Z}}_\mathcal{C}$
with $\mathcal{G}^\bullet$ bounded above complex of flat
$\underline{\mathbf{Z}}_\mathcal{C}$-modules is cofinal in the system.
(This follows from Modules on Sites, Lemma
\ref{sites-modules-lemma-module-quotient-flat} and
Derived Categories, Lemma \ref{derived-lemma-subcategory-left-resolution};
see discussion in Section \ref{section-flat}.)
Hence we can construct an inverse to the map
$H^0(\mathcal{C}, K) \longrightarrow H^0(\mathcal{C}, K_{ab})$
by representing an element $\xi \in H^0(\mathcal{C}, K_{ab})$ by a pair
$$
(s : \mathcal{G}^\bullet \to \underline{\mathbf{Z}}_\mathcal{C},
a : \mathcal{G}^\bullet \to \mathcal{K}^\bullet)
$$
with $\mathcal{G}^\bullet$ a bounded above complex of flat
$\underline{\mathbf{Z}}_\mathcal{C}$-modules and sending this to
$$
(\mathcal{G}^\bullet \otimes_{\underline{\mathbf{Z}}_\mathcal{C}}
\mathcal{O}_\mathcal{C}
\to \mathcal{O}_\mathcal{C},
\mathcal{G}^\bullet  \otimes_{\underline{\mathbf{Z}}_\mathcal{C}}
\mathcal{O}_\mathcal{C}
\to \mathcal{K}^\bullet)
$$
The only thing to note here is that the first arrow
is a quasi-isomorphism by
Lemmas \ref{lemma-derived-tor-quasi-isomorphism-other-side} and
\ref{lemma-bounded-flat-K-flat}.
We omit the detailed verification that this construction
is indeed an inverse.
\end{proof}

\end{document}
